% =====================================================================
% SPEC 2026 — IEEE Southern Power Electronics Conference
% Submission deadline: 31 July 2026
% Template: IEEEtran v1.8b+ (conference class)
% Compile: ./build.sh   (pdflatex → bibtex → pdflatex × 2)
% =====================================================================
\documentclass[conference, a4paper]{IEEEtran}

% --- Standard IEEE packages ---
\IEEEoverridecommandlockouts
\usepackage{cite}
\usepackage{amsmath,amssymb,amsfonts}
% \usepackage{algorithm,algpseudocode}  % uncomment after `tlmgr install algorithms`
% \usepackage{booktabs}                  % uncomment after `tlmgr install booktabs`
\usepackage{graphicx}
\usepackage{textcomp}
\usepackage{xcolor}
\usepackage{hyperref}
\hypersetup{
  colorlinks=true,
  linkcolor=black,
  citecolor=black,
  urlcolor=blue,
  pdftitle={Sort-and-Select Capacitor Balancing in a Single-Phase Modular Multilevel Converter},
  pdfauthor={Luiz Carlos Gili}
}

\def\BibTeX{{\rm B\kern-.05em{\sc i\kern-.025em b}\kern-.08em
    T\kern-.1667em\lower.7ex\hbox{E}\kern-.125emX}}

\begin{document}

% =====================================================================
% Title block
% =====================================================================
\title{
  Sort-and-Select Capacitor Balancing in a Single-Phase Modular
  Multilevel Converter: An Open-Source Reference Implementation
  in the Pulsim Simulator
}

\author{
  \IEEEauthorblockN{Luiz Carlos Gili}
  \IEEEauthorblockA{
    Independent Researcher \\
    Brazil \\
    \texttt{luizcarlosgili@gmail.com} \\
    ORCID: \href{https://orcid.org/0000-0002-5749-7199}{0000-0002-5749-7199}
  }
}

\maketitle

% =====================================================================
% Abstract
% =====================================================================
\begin{abstract}
The modular multilevel converter (MMC) is the dominant topology for
high-voltage DC transmission and a growing share of medium-voltage
drives. Its defining control challenge is the active balancing of
the floating sub-module (SM) capacitor voltages — without it, the
voltages diverge within a few line cycles and the multilevel
staircase output collapses. This paper presents an end-to-end
open-source reference implementation of the canonical
sort-and-select balancing algorithm in the Pulsim circuit simulator,
applied to a single-phase MMC with $N = 3$ SMs per arm under
phase-shifted-carrier (PSC) pulse-width modulation. Side-by-side
simulation results contrast the open-loop case — where the six
SM cap voltages diverge by more than 150~V within two line cycles
— with the closed-loop sort-and-select case, where all six
voltages remain locked within 0.03~V of the nominal
$V_C = V_{dc}/N$ value. The closed-loop fundamental
line-to-neutral peak matches the analytical $m_a V_{dc}/2$
prediction within $0.8\,\%$. The full simulation source,
including the closed-loop step-observer wiring of the
sort-and-select rule into the Pulsim solver, is available under
the MIT license at \texttt{github.com/lgili/Pulsim}.
\end{abstract}

\begin{IEEEkeywords}
Modular multilevel converter (MMC), capacitor voltage balancing,
sort-and-select, phase-shifted carrier PWM, piecewise-linear
simulation, open-source software, Pulsim.
\end{IEEEkeywords}

% =====================================================================
% I. Introduction
% =====================================================================
\section{Introduction}
\label{sec:intro}

% TODO §I (≈500 words, 0.75 page)
% 1. Why MMC matters (HVDC, MV drives, modularity).
% 2. The capacitor-balancing problem (1-2 sentences only — it's the
%    headline; details go in §III).
% 3. Why open-source matters: lock-in of commercial tools (PLECS,
%    PSIM, Simscape, Saber), and the resulting reproducibility
%    cost in published research.
% 4. Contribution: this paper presents the first MIT-licensed
%    end-to-end MMC + sort-and-select reference implementation
%    paired with an executable simulator.
% 5. Roadmap: §II builds the topology + modulator, §III states the
%    sort-and-select algorithm, §IV outlines the Pulsim
%    implementation, §V presents simulation results, §VI concludes.

\textit{Placeholder — to be written.}

% =====================================================================
% II. MMC Topology and Phase-Shifted Carrier PWM
% =====================================================================
\section{MMC Topology and Phase-Shifted Carrier PWM}
\label{sec:topology}

% TODO §II (≈700 words, 1 page)
% Subsections:
%   A. Half-bridge sub-module (with Fig 1 = SM schematic)
%      - 2-switch HB-SM diagram
%      - State table: INSERT (S1 ON, S2 OFF) -> v_SM = +V_C
%                     BYPASS (S1 OFF, S2 ON) -> v_SM = 0
%      - Forbidden states (cap-short, arm-open)
%   B. Arm structure
%      - v_arm = sum_i s_i v_C_i, takes N+1 discrete values
%      - Steady-state V_C = V_dc/N
%   C. PSC-PWM modulation
%      - N carriers phase-shifted by 2*pi/N
%      - n_insert(t) = count of carriers below the arm reference
%      - Effective load-side ripple at N * f_carrier
%      - Fig 2: 3 phase-shifted carriers + reference + n_insert output

\textit{Placeholder — to be written.}

% =====================================================================
% III. Sort-and-Select Balancing Algorithm
% =====================================================================
\section{The Sort-and-Select Balancing Algorithm}
\label{sec:sas}

% TODO §III (≈600 words, 1 page)
% 1. Why open-loop fails: naive modulator always inserts SM_0, SM_1, ...
%    => SM_0 hogs all charge, SM_{N-1} collapses (see Fig 4 in §V).
% 2. The sort-and-select rule (pseudocode in algorithm 1):
%
%    Input: n_insert (from PSC-PWM), V_C[1..N] (current cap voltages),
%           i_arm (current arm current sign)
%    if i_arm > 0:     // arm charging the inserted SMs
%        select the n_insert SMs with LOWEST V_C
%    else:             // arm discharging
%        select the n_insert SMs with HIGHEST V_C
%
% 3. Complexity: O(N log N) per switching update.
% 4. Why it works: at each instant, the algorithm picks the SMs
%    that will MOVE V_C toward the average. Provably stable for any
%    bounded load current.
% 5. Brief literature pointer (Saeedifard 2010; Marquardt 2018)
%    for theoretical analysis we don't repeat.

% TODO: when `algorithm,algpseudocode` are installed, replace this
% verbatim block with the proper \begin{algorithm}...\end{algorithm}
% formatting. Until then, the algorithm is described in prose +
% pseudocode comments below.
\textit{Algorithm pseudocode placeholder} — once `algorithm,
algpseudocode` are installed via TeX Live, replace this with the
properly formatted \texttt{Algorithm 1} block. For now, the rule is:

\begin{quote}\small\ttfamily
\textbf{Inputs:} $n_{\text{insert}}$ from PSC-PWM,
$V_C[1..N]$ cap voltages, $i_{\text{arm}}$ sign. \\
sort $\sigma$ s.t. $V_C[\sigma(1)] \le \cdots \le V_C[\sigma(N)]$. \\
\textbf{if} $i_{\text{arm}} > 0$
  \quad INSERT $\sigma(1), \ldots, \sigma(n_{\text{insert}})$ \\
\textbf{else}
  \quad INSERT $\sigma(N - n_{\text{insert}} + 1), \ldots, \sigma(N)$ \\
BYPASS the remaining SMs.
\end{quote}

\textit{Placeholder text around the algorithm — to be written.}

% =====================================================================
% IV. Implementation in Pulsim
% =====================================================================
\section{Implementation in Pulsim}
\label{sec:pulsim}

% TODO §IV (≈500 words, 0.75 page)
% 1. Pulsim brief: C++23 kernel + Python API, piecewise-linear (PWL)
%    state-space cache. Cite JOSS paper.
% 2. The CircuitBuilder code that wires up the MMC topology
%    (12 switches + 6 caps + 2 arm inductors + RL load). One short
%    snippet (≤ 15 lines).
% 3. The step_observer pattern: a closure shared between the
%    switch_fn (called every dt) and the observer (called every dt
%    AFTER the integration step). Observer reads V_C and i_arm out
%    of the state vector, sort-and-select runs in switch_fn at the
%    NEXT step.
% 4. Why this matters for reproducibility: same C++ kernel runs both
%    open-loop and closed-loop variants — only the switch_fn closure
%    differs. No recompilation, no separate binaries.

\textit{Placeholder — to be written.}

% =====================================================================
% V. Simulation Results
% =====================================================================
\section{Simulation Results}
\label{sec:results}

% TODO §V (≈800 words, 1.75 pages — the figure-heavy section)
% Sub-sections:
%   A. Operating point + simulation parameters (1 table)
%      V_dc=400V, N=3, V_C nom=133.3V, f_o=60Hz, f_carrier=1kHz,
%      RL load 24Ω + 5mH, single-phase, m_a=0.78
%   B. Open-loop: caps diverge (Fig 4)
%      Quote: "within 33 ms (2 line cycles), V_C[SM_u0] reaches
%      167.8 V (+25.8%), while V_C[SM_u2] collapses to 2.0 V."
%   C. Sort-and-select: caps locked (Fig 5)
%      Quote: "After the same 33 ms window with sort-and-select
%      active, the 6 cap voltages remain within 0.03 V of each
%      other."
%   D. 4-level arm voltage + 4-level AC output (Fig 6)
%      Quantitative: peak v_ac = 156.9 V, target = 155.6 V (0.8% err)
%   E. (Optional, if space) Output THD comparison with a 2-level VSI
%      driven at the same fundamental — quote NPC project result for
%      ballpark numbers.

\textit{Placeholder — to be written.}

% =====================================================================
% VI. Conclusion
% =====================================================================
\section{Conclusion}
\label{sec:conclusion}

% TODO §VI (≈200 words, 0.25 page)
% 1. Recap: open-source MMC + sort-and-select reference impl works,
%    cap-voltage error ≤ 0.03 V, fundamental ≤ 1% off analytical.
% 2. Why it matters: lowers the entry barrier for MMC research +
%    teaching.
% 3. Future work: 3-phase extension (~3× state space, beyond Pulsim
%    PWL cache today), circulating-current resonant suppression at
%    2*f_o, hardware-in-the-loop validation.

\textit{Placeholder — to be written.}

% =====================================================================
% Acknowledgement (optional, IEEE conf often omits)
% =====================================================================
% \section*{Acknowledgement}
% No external funding to acknowledge for this work.

% =====================================================================
% References
% =====================================================================
\bibliographystyle{IEEEtran}
\bibliography{refs}

\end{document}
